% -------------------------------------------------
% فصل 5
% -------------------------------------------------

\projectchapter{نتیجه‌گیری و پیشنهادها}{
در این فصل، دستاوردهای نهایی پروژه شناسایی زبان در داده‌های متنی و گفتاری جمع‌بندی می‌شوند. نتایج نشان دادند که خطوط لوله نظارت‌شده در محدوده مجموعه‌داده‌های پروژه عملکرد بسیار مناسبی دارند؛ بااین‌حال، نتیجه خوشه‌بندی به نوع داده و معیار انتخاب مدل وابسته است و جدایی طبیعی نمونه‌ها همیشه با برچسب زبان مطابقت ندارد. همچنین محدودیت‌های داده و پیاده‌سازی بیان می‌شوند و مسیرهای قابل بررسی برای توسعه آینده پیشنهاد می‌گردند.
}


% =================================================
\section{جمع‌بندی پروژه}
% =================================================

هدف این پروژه طراحی و ارزیابی یک سامانه شناسایی زبان برای دو نوع داده مستقل، یعنی متن و گفتار، بود. این دو مجموعه‌داده از نظر تعداد نمونه، زبان‌ها، شناسه‌ها و ساختار فایل با یکدیگر یکسان نیستند؛ ازاین‌رو برای هرکدام خط لوله جداگانه‌ای ایجاد شد و نتایج آن‌ها به‌عنوان یک آزمایش جفت‌شده یا مقایسه کنترل‌شده چندوجهی تفسیر نشدند.

در مسیر متن، داده‌های خام کشف و اعتبارسنجی شدند و پس از پاک‌سازی فاصله‌های اضافی، ویژگی‌های کاراکتری و واژگانی
\lr{TF-IDF}
استخراج شدند. سپس انتخاب 50 ویژگی برتر با معیار
\lr{Chi-Square}
درون خط لوله مدل انجام شد. قرارگرفتن بردارسازی و انتخاب ویژگی درون خط لوله باعث شد این مراحل در هر دور اعتبارسنجی فقط روی داده آموزشی برازش شوند و اطلاعات مجموعه آزمون وارد فرایند آموزش نشود.

در مسیر گفتار، ابتدا فایل‌های صوتی توسط نگارنده پاک‌سازی و تحلیل شدند و خروجی استخراج ویژگی در دو فایل
\lr{features.csv}
و
\lr{metadata.csv}
ذخیره شد. خط لوله آزمایش نهایی این دو فایل را به‌صورت فقط‌خواندنی دریافت کرد و آموزش مدل‌ها را از جدول 133 ویژگی عددی آغاز نمود. برای مدیریت مقادیر گمشده و تفاوت مقیاس ویژگی‌ها، جایگزینی مقادیر و استانداردسازی درون خطوط لوله طبقه‌بندی قرار گرفتند.

برای کاهش نشت اطلاعات، تقسیم آموزش و آزمون در هر دو مسیر به‌صورت گروه‌آگاه انجام شد. تنظیم ابرپارامترها و انتخاب طبقه‌بند فقط بر اساس میانگین
\lr{Macro F1}
در اعتبارسنجی متقاطع بخش آموزشی صورت گرفت و مجموعه آزمون نهایی در انتخاب مدل نقشی نداشت. در بخش خوشه‌بندی نیز برچسب زبان هنگام برازش، کاهش بُعد و انتخاب پارامترها استفاده نشد و معیارهای خارجی تنها پس از تشکیل خوشه‌ها محاسبه شدند.

در مجموع، برای هر یک از دو نوع داده دقیقاً چهار خانواده طبقه‌بندی و چهار خانواده خوشه‌بندی اجرا شد. بنابراین 16 نتیجه مربوط به خانواده‌های مدل تولید گردید. هفت دفترچه پروژه نیز به‌طور کامل اجرا شدند، مدل منتخب طبقه‌بندی هر مسیر ذخیره شد و آزمون استنتاج هر دو بسته مدل با موفقیت انجام گرفت. وضعیت اجرای کامل پروژه در فایل کنترل کیفیت برابر
\lr{passed}
ثبت شده است.


% =================================================
\section{مهم‌ترین نتایج}
% =================================================

مهم‌ترین یافته‌های حاصل از اجرای نهایی پروژه به شرح زیر هستند:

\begin{itemize}
	\item \textbf{طبقه‌بندی متن:}
	مدل
	\lr{Logistic Regression}
	بر اساس میانگین
	\lr{Macro F1}
	برابر با 0.9906 در اعتبارسنجی متقاطع گروه‌آگاه انتخاب شد. این مدل روی 211 نمونه آزمون نهایی به
	\lr{Accuracy}
	برابر با 0.9953 و
	\lr{Macro F1}
	برابر با 0.9946 رسید و فقط یک متن اسپانیایی را پرتغالی تشخیص داد. نتایج نشان می‌دهند که ترکیب ویژگی‌های کاراکتری و واژگانی برای زبان‌های موجود در این مجموعه‌داده مؤثر بوده است.

	\item \textbf{مقایسه طبقه‌بندهای متن:}
	مدل‌های
	\lr{KNN}،
	\lr{Decision Tree}
	و
	\lr{Multinomial Naive Bayes}
	نیز عملکرد مناسبی داشتند، اما امتیاز اعتبارسنجی آن‌ها از مدل منتخب کمتر بود. دقت آموزشی کامل در درخت تصمیم نشانه‌ای از ظرفیت بالای این مدل برای بیش‌برازش بود. همچنین
	\lr{Multinomial Naive Bayes}
	بیشترین تغییرپذیری میان دورهای اعتبارسنجی را نشان داد. در نتیجه، انتخاب
	\lr{Logistic Regression}
	فقط بر دقت آزمون متکی نبود و بر اساس معیار از پیش تعیین‌شده آموزشی انجام شد.

	\item \textbf{طبقه‌بندی گفتار:}
	مدل
	\lr{KNN}
	با میانگین
	\lr{Macro F1}
	برابر با 0.9662 در اعتبارسنجی متقاطع گروه‌آگاه انتخاب شد. این مدل روی 148 نمونه آزمون نهایی به
	\lr{Accuracy}
	برابر با 0.9932 و
	\lr{Macro F1}
	برابر با 0.9914 رسید و تنها یک نمونه آلمانی را کره‌ای پیش‌بینی کرد.

	\item \textbf{مقایسه طبقه‌بندهای گفتار:}
	سه مدل
	\lr{Logistic Regression}،
	\lr{Random Forest}
	و
	\lr{RBF SVM}
	روی مجموعه آزمون به دقت 1.0000 رسیدند. بااین‌حال، این نتیجه برای انتخاب مدل استفاده نشد؛ زیرا آزمون نهایی باید خارج از فرایند انتخاب باقی می‌ماند. اختلاف امتیاز اعتبارسنجی
	\lr{KNN}
	و
	\lr{Logistic Regression}
	حدود 0.0006 و بسیار کوچک‌تر از انحراف معیار اعتبارسنجی آن‌ها بود. بنابراین نمی‌توان برتری قطعی و عمومی
	\lr{KNN}
	را ادعا کرد؛ این مدل صرفاً مطابق قاعده از پیش تعیین‌شده پروژه انتخاب شده است.

	\item \textbf{خوشه‌بندی متن:}
	مدل‌های
	\lr{K-Means}
	و
	\lr{Agglomerative Clustering}
	هر شش زبان را در این مجموعه‌داده به‌طور کامل بازیابی کردند و برای هر سه معیار
	\lr{ARI}،
	\lr{NMI}
	و
	\lr{Purity}
	مقدار 1.0000 به دست آوردند. با وجود این، مدل
	\lr{DBSCAN}
	به دلیل بیشترین معیار داخلی
	\lr{Silhouette}
	برابر با 0.4265 به‌عنوان مدل منتخب بدون برچسب ثبت شد. این مدل 13 خوشه تشکیل داد و 143 نمونه، معادل 23.52 درصد داده‌ها، را نویز تشخیص داد. این نتیجه نشان می‌دهد که بهترین معیار هندسی لزوماً به معنای بهترین بازیابی برچسب معنایی زبان نیست.

	\item \textbf{خوشه‌بندی گفتار:}
	مدل
	\lr{Agglomerative Clustering}
	با مقدار
	\lr{Silhouette}
	برابر با 0.4709 طبق معیار داخلی انتخاب شد، اما یک خوشه شامل 719 نمونه و خوشه دیگر شامل تنها یک نمونه بود. مقدار
	\lr{ARI}
	آن برابر صفر و
	\lr{NMI}
	برابر با 0.0028 بود؛ بنابراین مقدار بالای
	\lr{Silhouette}
	در این حالت حاصل جداسازی یک نمونه دورافتاده است و به معنای کشف زبان‌ها نیست. در میان مدل‌های گفتار،
	\lr{K-Means}
	بیشترین انطباق خارجی را با
	\lr{ARI}
	برابر با 0.4983،
	\lr{NMI}
	برابر با 0.6731 و
	\lr{Purity}
	برابر با 0.9028 ایجاد کرد، ولی داده‌ها را به هشت خوشه و نه چهار خوشه متناظر با زبان‌ها تقسیم نمود.
\end{itemize}

در مجموع، نتایج طبقه‌بندی نشان دادند که زبان نمونه‌های موجود با استفاده از برچسب‌های آموزشی به‌خوبی قابل تشخیص است. در مقابل، نتایج خوشه‌بندی ثابت کردند که ساختار غالب بدون برچسب همیشه همان مفهوم مورد نظر پژوهشگر نیست. این تفاوت در گفتار آشکارتر بود؛ زیرا ویژگی‌های آکوستیکی ممکن است علاوه بر زبان، تحت تأثیر ویژگی‌های گوینده، منبع ضبط، مدت فایل و شرایط صوتی نیز قرار گیرند. بررسی دقیق سهم هر یک از این عوامل در آزمایش فعلی انجام نشده است و این موارد فقط توضیح‌های احتمالی محسوب می‌شوند.


% =================================================
\section{محدودیت‌های پروژه}
% =================================================

با وجود عملکرد مناسب مدل‌ها، نتایج باید با توجه به محدودیت‌های زیر تفسیر شوند:

\begin{itemize}
	\item مجموعه‌داده متن شامل 608 نمونه و شش زبان و مجموعه‌داده گفتار شامل 720 نمونه و چهار زبان است. این اندازه و تنوع برای ارائه یک خط پایه دانشگاهی مناسب است، اما برای نتیجه‌گیری عمومی درباره همه زبان‌ها، گویندگان و شرایط واقعی کافی نیست.

	\item داده‌های متن و گفتار مستقل هستند و زبان‌ها و نمونه‌های کاملاً یکسانی ندارند. بنابراین نتایج دو مسیر نباید به‌عنوان مقایسه مستقیم کیفیت متن و گفتار یا به‌عنوان یک سامانه چندوجهی جفت‌شده تفسیر شوند.

	\item دقت بسیار بالای متن ممکن است تا حدی ناشی از تفاوت خط، الگوهای واژگانی مشخص و ویژگی‌های منبع داده باشد. تعمیم این عملکرد به متن‌های بسیار کوتاه، دارای غلط املایی، نویسه‌گردانی‌شده، آمیخته با چند زبان یا متعلق به منابع دیگر هنوز ارزیابی نشده است.

	\item ویژگی‌های گفتار به شکل میانگین و انحراف معیار ویژگی‌های کل فایل ذخیره شده‌اند. این نمایش، ترتیب زمانی و تغییرات محلی سیگنال را حفظ نمی‌کند و ممکن است بخشی از اطلاعات زبانی گفتار را از بین ببرد.

	\item در داده گفتار یک نمونه با مدت تقریبی 1534 ثانیه وجود دارد. این نمونه حذف نشده است و تحلیل حساسیت با حذف آن یا کنارگذاشتن ویژگی مدت زمان انجام نشده؛ در نتیجه میزان اثر دقیق آن بر مدل‌ها مشخص نیست.

	\item گروه‌ها از الگوی نام فایل استخراج شده‌اند. این روش برای کاهش هم‌پوشانی منبع میان آموزش و آزمون مفید است، اما در همه موارد معادل شناسه تأییدشده گوینده نیست. بنابراین احتمال باقی‌ماندن برخی وابستگی‌های ناشناخته میان نمونه‌ها کاملاً منتفی نیست.

	\item برای ارزیابی نهایی از یک تقسیم گروه‌آگاه آموزش و آزمون استفاده شده است. اعتبارسنجی متقاطع روی بخش آموزشی عدم قطعیت انتخاب مدل را تا حدی نشان می‌دهد، اما ارزیابی روی مجموعه‌داده خارجی و مستقل همچنان لازم است.

	\item انتخاب مدل خوشه‌بندی بر اساس یک معیار داخلی می‌تواند در حضور نویز، خوشه‌های بسیار کوچک یا نمونه‌های دورافتاده گمراه‌کننده باشد. نمونه روشن آن، مدل سلسله‌مراتبی گفتار است. پایداری خوشه‌ها در تقسیم‌ها یا نمونه‌گیری‌های مجدد نیز در این پروژه بررسی نشده است.

	\item بسته استنتاج متن، متن خام را مستقیماً دریافت می‌کند؛ اما بسته فعلی گفتار به یک ردیف شامل همان 133 ویژگی استخراج‌شده نیاز دارد و فایل صوتی خام را مستقیماً نمی‌پذیرد.

	\item مسئله تشخیص زبان ناشناخته، تنظیم اطمینان پیش‌بینی، گفتار یا متن چندزبانه و تغییر توزیع داده در محیط واقعی در محدوده این نسخه قرار نداشته‌اند.
\end{itemize}

این محدودیت‌ها اعتبار نتایج اجراشده را از بین نمی‌برند، بلکه دامنه صحیح تفسیر آن‌ها را مشخص می‌کنند. اعداد گزارش‌شده بیانگر عملکرد مدل‌ها روی تقسیم‌ها و داده‌های همین پروژه هستند و نباید بدون آزمایش تکمیلی به مجموعه‌داده‌های دیگر تعمیم داده شوند.


% =================================================
\section{پیشنهاد برای توسعه آینده}
% =================================================

با توجه به نتایج و محدودیت‌های مشاهده‌شده، پیشنهادهای توسعه آینده به ترتیب اولویت عبارت‌اند از:

\begin{enumerate}
	\item \textbf{ارزیابی روی داده‌های بزرگ‌تر و مستقل:}
	مهم‌ترین گام، آزمایش مدل‌های ذخیره‌شده روی منابع جدید با گویندگان، موضوع‌ها، لهجه‌ها، ابزارهای ضبط و شرایط نویزی متنوع است. افزودن زبان‌ها و نمونه‌های کوتاه یا دشوار نیز امکان بررسی واقع‌بینانه‌تر تعمیم‌پذیری را فراهم می‌کند.

	\item \textbf{بهبود فراداده و تقسیم گفتار:}
	استفاده از شناسه تأییدشده گوینده و منبع ضبط، تقسیم دقیق‌تر گروه‌آگاه را ممکن می‌سازد. همچنین بهتر است یک مجموعه ارزیابی کاملاً مستقل از گوینده و منبع آموزشی نگهداری شود.

	\item \textbf{تکمیل مسیر استنتاج صوت خام:}
	می‌توان مراحل پاک‌سازی و استخراج ویژگی موجود را به‌صورت یک مسیر اختیاری و بازتولیدپذیر به بسته استنتاج متصل کرد تا کاربر به جای جدول ویژگی، فایل صوتی را مستقیماً وارد کند. در این حالت باید ترتیب، نام و تعریف دقیق هر 133 ویژگی با قرارداد مدل ذخیره‌شده یکسان باقی بماند.

	\item \textbf{تحلیل حساسیت ویژگی‌های گفتار:}
	اثر نمونه بسیار طولانی، ویژگی مدت زمان و روش خلاصه‌سازی کل فایل باید با آزمایش‌های کنترل‌شده بررسی شود. قطعه‌بندی فایل‌های طولانی و مقایسه نتیجه با و بدون نمونه‌های دورافتاده می‌تواند علت رفتار خوشه‌بندی را روشن‌تر کند.

	\item \textbf{بررسی نمایش‌های زمانی گفتار:}
	در کنار خط پایه کلاسیک، می‌توان دنباله‌های
	\lr{MFCC}،
	مدل‌های
	\lr{CNN/RNN}
	یا تعبیه‌های ازپیش‌آموزش‌دیده گفتار را ارزیابی کرد. این مدل‌ها ممکن است اطلاعات زمانی بیشتری حفظ کنند، اما هزینه محاسباتی و نیاز داده‌ای آن‌ها باید با خط پایه قابل اجرای پردازنده مرکزی مقایسه شود.

	\item \textbf{آزمایش‌های مقاوم‌بودن متن:}
	مقایسه جداگانه ویژگی‌های کاراکتری و واژگانی، تغییر بازه
	\lr{N-Gram}،
	تغییر تعداد ویژگی‌های منتخب و افزودن نویز مصنوعی مانند غلط املایی یا نویسه‌گردانی می‌تواند وابستگی مدل به نشانه‌های سطحی را مشخص کند.

	\item \textbf{تشخیص زبان ناشناخته و کالیبراسیون:}
	افزودن آستانه اطمینان، کالیبراسیون احتمال و سازوکاری برای رد نمونه‌هایی که متعلق به زبان‌های آموزش‌دیده نیستند، سامانه را برای استفاده عملی ایمن‌تر می‌کند.

	\item \textbf{ارزیابی پایداری خوشه‌بندی:}
	بهتر است در کنار معیارهای داخلی و خارجی، حداقل اندازه خوشه، پایداری در نمونه‌گیری مجدد و حساسیت به پارامترها نیز سنجیده شود. روش‌های مقاوم‌تر چگالی‌محور مانند
	\lr{HDBSCAN}
	و تحلیل ترکیب خوشه‌ها بر اساس گوینده یا منبع، پس از پایان برازش بدون برچسب، می‌توانند درک بهتری از ساختار داده ارائه دهند.

	\item \textbf{بررسی سامانه چندوجهی در صورت تهیه داده جفت‌شده:}
	اگر در آینده برای هر نمونه صوتی، متن متناظر و برچسب‌های مشترک فراهم شود، می‌توان روش‌های ترکیب متن و گفتار را بررسی کرد. داده‌های مستقل فعلی برای نتیجه‌گیری معتبر درباره ادغام چندوجهی مناسب نیستند.
\end{enumerate}

در نهایت، پروژه حاضر یک خط پایه کامل، بازتولیدپذیر و قابل اجرا روی پردازنده مرکزی برای شناسایی زبان فراهم کرده است. مهم‌ترین دستاورد آن تنها دستیابی به دقت بالا نیست، بلکه رعایت جداسازی داده آزمون، برازش ویژگی‌ها درون خط لوله، انتخاب مدل بر اساس اعتبارسنجی و گزارش شفاف تفاوت میان معیارهای داخلی و خارجی خوشه‌بندی است. بر اساس شواهد فعلی، طبقه‌بندی نظارت‌شده مناسب‌ترین مسیر برای کاربرد اصلی این پروژه است؛ درحالی‌که خوشه‌بندی بیشتر به‌عنوان ابزاری برای کشف و تحلیل ساختار پنهان داده‌ها ارزش دارد. توسعه آینده باید پیش از افزایش پیچیدگی مدل، بر تنوع داده، اعتبارسنجی خارجی و مقاوم‌بودن سامانه تمرکز کند.

